

@article{feynman1982simulating,
  added-at = {2019-02-26T21:36:34.000+0100},
  author = {Feynman, Richard P},
  biburl = {https://www.bibsonomy.org/bibtex/213ff745d56ca88760b0e36f2b112e7a4/rspreeuw},
  interhash = {e7b5cbbbe625b40aca61c24bdbee745d},
  intrahash = {13ff745d56ca88760b0e36f2b112e7a4},
  journal = {International journal of theoretical physics},
  keywords = {imported},
  number = {6/7},
  pages = {467--488},
  publisher = {World Scientific},
  timestamp = {2019-02-26T21:36:34.000+0100},
  title = {Simulating physics with computers},
  volume = 21,
  year = 1982
}



@article{
doi:10.1126/science.273.5278.1073,
author = {Seth Lloyd },
title = {Universal Quantum Simulators},
journal = {Science},
volume = {273},
number = {5278},
pages = {1073-1078},
year = {1996},
doi = {10.1126/science.273.5278.1073},
URL = {https://www.science.org/doi/abs/10.1126/science.273.5278.1073},
eprint = {https://www.science.org/doi/pdf/10.1126/science.273.5278.1073},
abstract = {Feynman's 1982 conjecture, that quantum computers can be programmed to simulate any local quantum system, is shown to be correct.}
}



@article{PhysRevLett.114.090502,
  title = {Simulating Hamiltonian Dynamics with a Truncated Taylor Series},
  author = {Berry, Dominic W. and Childs, Andrew M. and Cleve, Richard and Kothari, Robin and Somma, Rolando D.},
  journal = {Phys. Rev. Lett.},
  volume = {114},
  issue = {9},
  pages = {090502},
  numpages = {5},
  year = {2015},
  month = {Mar},
  publisher = {American Physical Society},
  doi = {10.1103/PhysRevLett.114.090502},
  url = {https://link.aps.org/doi/10.1103/PhysRevLett.114.090502}
}

@article{Berry2015HamiltonianSW,
  title={Hamiltonian Simulation with Nearly Optimal Dependence on all Parameters},
  author={Dominic W. Berry and Andrew M. Childs and Robin Kothari},
  journal={2015 IEEE 56th Annual Symposium on Foundations of Computer Science},
  year={2015},
  pages={792-809}
}


@article{PhysRevLett.118.010501,
  title = {Optimal Hamiltonian Simulation by Quantum Signal Processing},
  author = {Low, Guang Hao and Chuang, Isaac L.},
  journal = {Phys. Rev. Lett.},
  volume = {118},
  issue = {1},
  pages = {010501},
  numpages = {5},
  year = {2017},
  month = {Jan},
  publisher = {American Physical Society},
  doi = {10.1103/PhysRevLett.118.010501},
  url = {https://link.aps.org/doi/10.1103/PhysRevLett.118.010501}
}

@article{Low2019hamiltonian,
  doi = {10.22331/q-2019-07-12-163},
  url = {https://doi.org/10.22331/q-2019-07-12-163},
  title = {Hamiltonian {S}imulation by {Q}ubitization},
  author = {Low, Guang Hao and Chuang, Isaac L.},
  journal = {{Quantum}},
  issn = {2521-327X},
  publisher = {{Verein zur F{\"{o}}rderung des Open Access Publizierens in den Quantenwissenschaften}},
  volume = {3},
  pages = {163},
  month = jul,
  year = {2019}
}


@article{TCS-066,
url = {http://dx.doi.org/10.1561/0400000066},
year = {2015},
volume = {10},
journal = {Foundations and Trends® in Theoretical Computer Science},
title = {Quantum Hamiltonian Complexity},
doi = {10.1561/0400000066},
issn = {1551-305X},
number = {3},
pages = {159-282},
author = {Sevag Gharibian and Yichen Huang and Zeph Landau and Seung Woo Shin}
}


@article{Kitaev1995QuantumMA,
  title={Quantum measurements and the Abelian Stabilizer Problem},
  author={Alexei Y. Kitaev},
  journal={Electron. Colloquium Comput. Complex.},
  year={1995},
  volume={TR96}
}


@article{doi:10.1137/16M1087072,
author = {Childs, Andrew M. and Kothari, Robin and Somma, Rolando D.},
title = {Quantum Algorithm for Systems of Linear Equations with Exponentially Improved Dependence on Precision},
journal = {SIAM Journal on Computing},
volume = {46},
number = {6},
pages = {1920-1950},
year = {2017},
doi = {10.1137/16M1087072},

URL = { 
    
        https://doi.org/10.1137/16M1087072
    
    

},
eprint = { 
    
        https://doi.org/10.1137/16M1087072
    
    

}
,
    abstract = { Harrow, Hassidim, and Lloyd [Phys. Rev. Lett., 103 (2009), 150502] showed that for a suitably specified \$N \times N\$ matrix \$A\$ and an \$N\$-dimensional vector \$\vec{b}\$, there is a quantum algorithm that outputs a quantum state proportional to the solution of the linear system of equations \$A\vec{x} = \vec{b}\$. If \$A\$ is sparse and well-conditioned, their algorithm runs in time \${poly}(\log N, 1/\epsilon)\$, where \$\epsilon\$ is the desired precision in the output state. We improve this to an algorithm whose running time is polynomial in \$\log(1/\epsilon)\$, exponentially improving the dependence on precision while keeping essentially the same dependence on other parameters. Our algorithm is based on a general technique for implementing any operator with a suitable Fourier or Chebyshev series representation. This allows us to bypass the quantum phase estimation algorithm, whose dependence on \$\epsilon\$ is prohibitive. }
}

@article{osti_1623945,
title = {A variational eigenvalue solver on a photonic quantum processor},
author = {Peruzzo, Alberto and McClean, Jarrod and Shadbolt, Peter and Yung, Man-Hong and Zhou, Xiao-Qi and Love, Peter J. and Aspuru-Guzik, Alán and O’Brien, Jeremy L.},
abstractNote = {Quantum computers promise to efficiently solve important problems that are intractable on a conventional computer. For quantum systems, where the physical dimension grows exponentially, finding the eigenvalues of certain operators is one such intractable problem and remains a fundamental challenge. The quantum phase estimation algorithm efficiently finds the eigenvalue of a given eigenvector but requires fully coherent evolution. Here we present an alternative approach that greatly reduces the requirements for coherent evolution and combine this method with a new approach to state preparation based on ansa¨tze and classical optimization. We implement the algorithm by combining a highly reconfigurable photonic quantum processor with a conventional computer. We experimentally demonstrate the feasibility of this approach with an example from quantum chemistry—calculating the ground-state molecular energy for He–H þ . The proposed approach drastically reduces the coherence time requirements, enhancing the potential of quantum resources available today and in the near future.},
doi = {10.1038/ncomms5213},
journal = {Nature Communications},
number = 1,
volume = 5,
place = {United States},
year = {2014},
month = {7}
}


@article{Motta_2019,
	doi = {10.1038/s41567-019-0704-4},
  
	url = {https://doi.org/10.1038%2Fs41567-019-0704-4},
  
	year = 2019,
	month = {nov},
  
	publisher = {Springer Science and Business Media {LLC}
},
  
	volume = {16},
  
	number = {2},
  
	pages = {205--210},
  
	author = {Mario Motta and Chong Sun and Adrian T. K. Tan and Matthew J. O'Rourke and Erika Ye and Austin J. Minnich and Fernando G. S. L. Brand{\~{a}}o and Garnet Kin-Lic Chan},
  
	title = {Determining eigenstates and thermal states on a quantum computer using quantum imaginary time evolution},
  
	journal = {Nature Physics}}}


 @article{Parrish2019QuantumFD,
  title={Quantum Filter Diagonalization: Quantum Eigendecomposition without Full Quantum Phase Estimation},
  author={Robert M. Parrish and Peter Leonard McMahon},
  journal={arXiv: Quantum Physics},
  year={2019}
}

@article{Jansen_2007,
	doi = {10.1063/1.2798382},
  
	url = {https://doi.org/10.1063\%2F1.2798382},
  
	year = 2007,
	month = {oct},
  
	publisher = {{AIP} Publishing},
  
	volume = {48},
  
	number = {10},
  
	pages = {102111},
  
	author = {Sabine Jansen and Mary-Beth Ruskai and Ruedi Seiler},
  
	title = {Bounds for the adiabatic approximation with applications to quantum computation},
  
	journal = {Journal of Mathematical Physics}
}


@article{PhysRevA.77.012326,
  title = {Quantum computational method of finding the ground-state energy and expectation values},
  author = {Oh, Sangchul},
  journal = {Phys. Rev. A},
  volume = {77},
  issue = {1},
  pages = {012326},
  numpages = {5},
  year = {2008},
  month = {Jan},
  publisher = {American Physical Society},
  doi = {10.1103/PhysRevA.77.012326},
  url = {https://link.aps.org/doi/10.1103/PhysRevA.77.012326}
}


@article{ge2018faster, 
      author = {Yimin Ge and Jordi Tura and J. Ignacio Cirac},
      title = {Faster ground state preparation and high-precision ground energy estimation with fewer qubits},
      month = {12},
      year = {2018},
      eprint = {1712.03193},
      archivePrefix = {arXiv},
      primaryClass = {quant-ph}
}


@article{Lin_2020,
	doi = {10.22331/q-2020-12-14-372},
  
	url = {https://doi.org/10.22331%2Fq-2020-12-14-372},
  
	year = 2020,
	month = {dec},
  
	publisher = {Verein zur Forderung des Open Access Publizierens in den Quantenwissenschaften},
  
	volume = {4},
  
	pages = {372},
  
	author = {Lin Lin and Yu Tong},
  
	title = {Near-optimal ground state preparation},
  
	journal = {Quantum}
}}


@article{brown2021quantum, 
      author={Adam R. Brown and Hrant Gharibyan and Stefan Leichenauer and Henry W. Lin and Sepehr Nezami and Grant Salton and Leonard Susskind and Brian Swingle and Michael Walter},
      title={Quantum Gravity in the Lab: Teleportation by Size and Traversable Wormholes},
      year={2021},
      eprint={1911.06314},
      archivePrefix={arXiv},
      primaryClass={quant-ph}
}




@article{Maldacena_2003,
	doi = {10.1088/1126-6708/2003/04/021},
  
	url = {https://doi.org/10.1088%2F1126-6708%2F2003%2F04%2F021},
  
	year = 2003,
	month = {apr},
  
	publisher = {Springer Science and Business Media {LLC}
},
  
	volume = {2003},
  
	number = {04},
  
	pages = {021--021},
  
	author = {Juan Maldacena},
  
	title = {Eternal black holes in anti-de Sitter},
  
	journal = {Journal of High Energy Physics}
}}

@article{Cottrell:2018ash,
    author = "Cottrell, William and Freivogel, Ben and Hofman, Diego M. and Lokhande, Sagar F.",
    title = "{How to Build the Thermofield Double State}",
    eprint = "1811.11528",
    archivePrefix = "arXiv",
    primaryClass = "hep-th",
    doi = "10.1007/JHEP02(2019)058",
    journal = "JHEP",
    volume = "02",
    pages = "058",
    year = "2019"
}

@inproceedings{Gily_n_2019_bis,
	doi = {10.1145/3313276.3316366},
  
	url = {https://doi.org/10.1145\%2F3313276.3316366},
  
	year = 2019,
	month = {jun},
  
	publisher = {{ACM}
},
  
	author = {Andr{\'{a}}s Gily{\'{e}}n and Yuan Su and Guang Hao Low and Nathan Wiebe},
  
	title = {Quantum singular value transformation and beyond: exponential improvements for quantum matrix arithmetics},
  
	booktitle = {Proceedings of the 51st Annual {ACM} {SIGACT} Symposium on Theory of Computing}
}

@inbook{Gily_n_2019,
  author       = {Andr{\'{a}}s Gily{\'{e}}n and Yuan Su and Guang Hao Low and Nathan Wiebe},
  title        = {Quantum singular value transformation and beyond: exponential improvements for quantum matrix arithmetics},
  booktitle    = {Proceedings of the 51st Annual {ACM} {SIGACT} Symposium on Theory of Computing},
  publisher    = {{ACM}},
  year         = {2019},
  editor       = {},
  pages        = {},
  chapter      = {},
  address      = {},
  edition      = {},
  month        = {June},
  note         = {DOI:10.1145/3313276.3316366, URL:https://doi.org/10.1145\%2F3313276.3316366}
}
